% ============================================================================
\section{Sparse OPD Updates and Teacher Subnetwork Overlap}
\label{sec:subnetworks}
% ============================================================================

We investigate three questions in sequence: whether OPD updates are sparse,
whether different teachers update overlapping subnetworks of the same MOPD
student, and whether larger teacher distribution differences correspond to
larger subnetwork differences. The following experiments are planned;
their outcomes are not yet established.

\subsection{Do sparse OPD updates persist under multiple teachers?}

Sparse parameter changes have been reported for RL fine-tuning
\citep{mukherjee2025subnetworks} and OPD
\citep{yu2026denseupdates,shen2026opdgeometry}. We first reproduce this
measurement in our setting and extend the comparison to MOPD. From the
same initialization $\theta_0$, we measure single-teacher deltas
$\Delta_{i,t}^{\mathrm{single}}$ and the joint delta
$\Delta_t^{\mathrm{joint}}=\theta_t-\theta_0$, comparing training curves
separately against domain exposure and compute. For an $n$-coordinate delta $v$,
\begin{equation}
 s_\tau(v)=\frac1n\sum_{a=1}^n
       \mathbf 1\{|v_a|\leq\tau\}.
 \label{subnet_sparsity}
\end{equation}
We report threshold sweeps, update norms, and top-coordinate energy
concentration using FP32 master-weight deltas, with bf16 checkpoint
measurements for comparison. This distinguishes small accumulated changes
from rounding-induced zeros; endpoint sparsity does not establish that
instantaneous gradients are exactly sparse. This first experiment is a
replication and extension, not a claim to discover OPD sparsity.

\subsection{Which subnetworks do teachers update in the same student?}

At selected checkpoints $\theta_t$ of one MOPD run, each teacher processes
its routed student-rollout batch independently. Every diagnostic branch
starts from the identical student and optimizer state. We record each
teacher's gradient $g_{i,t}$ and virtual optimizer displacement $u_{i,t}$
without advancing the shared training run. These local probes describe
each teacher's update tendency at the actual joint-training checkpoint;
they do not decompose AdamW's joint update additively.

For either probe $v_{i,t}$, define its threshold support and pairwise
Jaccard overlap as
\begin{equation}
 A_{i,t}^{\tau}=\{a:|v_{i,t,a}|>\tau\},\qquad
 J_{ij,t}=\frac{|A_{i,t}\cap A_{j,t}|}{|A_{i,t}\cup A_{j,t}|}.
 \label{subnet_overlap}
\end{equation}
We also compare masks retaining the same top fraction of coordinates,
globally and within layers, against density-matched random masks.
Threshold masks reflect update scale as well as location; fixed-density
masks emphasize location. Raw-gradient masks check whether shared momentum
or decay dominates virtual-step overlap. Independent single-teacher
endpoint supports provide a secondary comparison, but cannot identify
teacher contributions along the MOPD trajectory. Signed alignment is an
auxiliary diagnostic: overlapping coordinates can support compatible or
opposing updates, so low overlap is not assumed preferable.

\subsection{Do more different teachers update more different subnetworks?}

At each probe checkpoint, all teachers score the same cached student-prefix
bank $\mathcal C_t$. With $q_i$ denoting teacher $i$'s distribution, measure
\begin{equation}
\begin{aligned}
 D_{ij,t}^{\mathrm{TT}}&=\mathbb E_{c\sim\mathcal C_t}
          \operatorname{JS}(q_i(\cdot\mid c),q_j(\cdot\mid c)),\\
 D_{i,t}^{\mathrm{TS}}&=\mathbb E_{c\sim\mathcal C_t}
          \operatorname{JS}(q_i(\cdot\mid c),p_{\theta_t}(\cdot\mid c)).
\end{aligned}
\label{subnet_distances}
\end{equation}
We plot teacher-pair divergence against subnetwork distance $1-J_{ij,t}$
at the same checkpoint, recording both teacher--student gaps as auxiliary
descriptors. A small common-context check recomputes all teachers' masks
on identical prefixes, assessing whether routed-domain differences explain
the observed association. Code and science currently share one Qwen3-4B
teacher; comparisons between these routes measure context differences,
not differences between independent teachers. We report pair trajectories
and run-level uncertainty without treating shared teachers or repeated
checkpoints as independent observations. A positive distance--difference
relationship is a hypothesis to test, not a consequence of sparsity or an
assumption about optimal MOPD behavior.
